\section{案例 v1.1（实现中）：让物理页与堆对象都拥有可回收生命周期}
\label{case:v11-physical-memory}

\begin{statebox}
\textbf{状态边界。} 项目当前发布基线仍是 v1.0；v1.1 已由 ADR 接受并实现三项增量。
ADR 0017 建立动态页帧元数据、高半区物理直映和 64 GiB QEMU 容量路径；
ADR 0020 把 64 KiB 高半区单调堆升级为可释放、可合并并能完整校验的通用堆。
ADR 0022 又以逐阶 Free/Allocated 双位图建立 buddy 连续块的分裂、合并和
精确释放，并通过 \(80/80\) 项完整回归。README、roadmap、架构文档和测试文档
已经同步接受该增量；但 CMake 工程元数据仍为 v0.11，type cache、KVA、
动态内核栈和页表回收也未闭合。因此 buddy 是\textbf{已接受的 v1.1 增量}，
却仍不能被解释成“v1.1 已经发布完成”。
\end{statebox}

v0.6 已经能解析 E820 内存图，却只管理 Stage 1 恒等映射覆盖的低 64 MiB。
这在当时是一条安全边界：页表代码可把低端物理地址直接当作同值指针使用，
页帧状态也只需一份固定数组。进入 v1.1 后，QEMU 的 E820 可以报告 64 GiB
可用 RAM；若内核仍输出 64 MiB 的 managed bytes，问题并不在“机器没有内存”，
而在软件把\emph{启动时能直接访问的范围}误当成了\emph{长期能够管理的范围}。

\subsection{四个上限回答的是四个不同问题}

决定物理内存容量时，至少要区分下列四个边界。把它们压成一个
\code{memory\_size}，会使越界、地址洞和页表覆盖错误无法定位。

\begin{longtable}{|L{0.21\linewidth}|L{0.25\linewidth}|L{0.28\linewidth}|L{0.16\linewidth}|}
\toprule
边界 & 它回答的问题 & v1.1 的处理方式 & 不能推出什么 \\
\midrule
\endhead
E820 type 1 区间 & 固件声明哪些物理区间可作为普通 RAM & 逐项验证排序、重叠、
长度与溢出，只把可用区间加入页帧账本 & 地址在最高端点以下不等于可分配
\\\tablerowrule
处理器物理地址宽度 & 页表项中的物理地址最多能表达多大 & 从
\code{CPUID.80000008H:EAX[7:0]} 取得位数并计算上限 & E820 中每段都一定合法
\\\tablerowrule
高半区直映窗口 & 内核预留的虚拟窗口最多能覆盖多少物理地址 & 四级分页下为
direct map 保留 64 TiB，并检查基址加法与窗口溢出 & 窗口内地址已经建立映射
\\\tablerowrule
Stage 1 恒等映射 & 切换新 CR3 以前，哪些物理页能以同值地址解引用 & 继续保留
低 64 MiB 作为启动期建表与元数据落点 & 内核长期只能管理低 64 MiB \\
\bottomrule
\end{longtable}

受管最高物理地址因此取“最高可用 RAM 端点、处理器物理地址上限、直映窗口容量”
三者的最小值，再向下按 4 KiB 对齐。这个值只是 PFN 账本的覆盖上限；每个
E820 洞仍保持 Unavailable，不能因为它落在 \([0,\text{managed limit})\)
内部就变成 Free。

\begin{pdfdiagram}
\node[os state, text width=3.0cm] (e820)
  {E820 可用区间\\RAM 端点与地址洞\\决定候选 PFN};
\node[os compute, text width=3.0cm, right=14mm of e820] (limit)
  {容量裁剪\\物理地址宽度\\直映窗口容量};
\node[os memory, text width=3.0cm, right=14mm of limit] (ledger)
  {页帧状态账本\\每 PFN 两位\\洞保持 Unavailable};
\node[os control, text width=3.0cm, below=15mm of e820] (identity)
  {低 64 MiB 恒等映射\\新 CR3 前可解引用\\限制启动分配};
\node[os memory, text width=3.0cm, right=14mm of identity] (direct)
  {高半区 direct map\\新 CR3 后可解引用\\覆盖 type 1 RAM};
\node[os state, text width=3.0cm, right=14mm of direct] (owner)
  {所有权状态\\Free/Allocated\\Reserved 与映射分离};
\draw[os strong bus] (e820) -- (limit);
\draw[os strong bus] (limit) -- (ledger);
\draw[os strong bus] (identity) -- node[os label, above] {启动访问} (direct);
\draw[os strong bus] (ledger) -- (owner);
\draw[os strong bus] (direct) -- node[os label, above] {访问路径} (owner);
\draw[os feedback] (limit.south) to[bend right=13]
  node[os label, below] {容量不由低端映射决定} (identity.north);
\end{pdfdiagram}
\diagramnote{容量、访问和所有权是三套状态。E820 与处理器能力决定哪些 PFN
可以进入账本；恒等映射和 direct map 决定内核怎样访问物理页；页帧状态才决定
某一页能否分配。任何一条箭头都不能替代另外两条。}

\subsection{两位页帧状态为什么仍可能需要四 MiB 以上}

页大小为 4096 字节，每字节容纳四个两位状态。若没有地址洞，64 GiB RAM
包含
\[
  \frac{64 \times 2^{30}}{2^{12}} = 2^{24}
\]
个页帧，最少需要 \(2^{24}/4=2^{22}\) 字节，即 4 MiB 状态存储。
但元数据下标是 PFN，不是“第几个可用页”。若 3 GiB 可用区之后有 1 GiB 洞，
高端 RAM 再延伸到 65 GiB 物理地址，那么状态数组必须覆盖洞后的 PFN，实际
尺寸会略大于 4 MiB；当前 64 GiB QEMU 主规格计算出的状态存储为
\code{0x410000} 字节。洞对应的状态保持 Unavailable。反过来，高端保留设备
窗口只要不包含可用 RAM，就不能仅凭自身端点把状态数组扩大到该窗口。

v0.6 的固定 4096 字节数组恰好只能表达低 64 MiB：
\[
  4096\ \text{B}\times 4\ \text{frames/B}\times 4096\ \text{B/frame}
  =64\ \text{MiB}.
\]
因此，仅把 \code{MANAGED\_BYTES} 常量改成 64 GiB 会让数组越界。v1.1 先根据
managed limit 计算元数据尺寸，再从低 64 MiB 的可用区间中寻找一段连续、
页对齐且不与固件、Stage 1、内核、BootInfo、启动栈及临时缓冲重叠的空间。
若找不到，启动必须明确失败，不能静默退回较小容量。

\begin{lstlisting}
prepare_physical_memory(e820, boot_objects):
    // 容量同时受固件事实、处理器可表达范围和直映窗口约束。
    cpu_limit = physical_limit_from_cpuid()
    map_limit = direct_map_capacity()
    managed_limit = align_down(min(highest_usable_end(e820),
                                   cpu_limit, map_limit), PAGE_SIZE)

    // 两位状态按 PFN 编号，地址洞也占据索引，但初始保持不可用。
    metadata_bytes = align_up(ceil(page_count(managed_limit) / 4),
                              PAGE_SIZE)
    metadata = find_low_identity_mapped_range(
        e820, boot_objects, metadata_bytes, PAGE_SIZE)
    if metadata does not exist:
        return FRAME_STATE_STORAGE_UNAVAILABLE

    // 先把全部状态置为不可用，再只开放合法 E820 type 1 页。
    initialize_frame_ledger(metadata, managed_limit, default=UNAVAILABLE)
    mark_type1_pages_free(e820)
    reserve_boot_objects(boot_objects, metadata)
    return SUCCEEDED
\end{lstlisting}

\subsection{direct map 提供访问路径，不授予页帧所有权}

新直映窗口起于 \code{0xFFFF888000000000}。对一个被纳入直映的物理地址
\(p\)，内核使用
\[
  v=\mathtt{0xFFFF888000000000}+p
\]
得到可解引用的高半区虚拟地址。页表项仍保存 \(p\)，CR3 仍要求物理地址；
只有 C++ 代码准备读写页表帧、清零用户页或复制 ELF 段时，才把物理地址转换成
当前地址空间中的虚拟地址。

这条规则消除了 v0.6 中隐含的
\code{reinterpret\_cast<void *>(physicalAddress)}。在旧 CR3 下，页表管理器
使用“虚拟基址为 0、只从低 64 MiB 分配页表帧”的访问策略；激活新 CR3 后，
同一管理器切换为“虚拟基址为 direct-map base、可访问全部受管物理页”的策略。
算法仍处理物理页号，解引用方式则随控制权阶段改变。

直映只覆盖完整的 E820 type 1 RAM，权限为 supervisor、read/write、NX 和普通
缓存。对齐且完整落在可用区间内部的部分优先使用 2 MiB 大页，区间边界与尾部
使用 4 KiB 页。保留区、未知区和 MMIO 不会因为落入窗口范围而自动变成普通
缓存 RAM；例如 LAPIC 仍由独立的 cache-disabled 映射处理。

\subsection{为什么换 CR3 必须是一笔启动事务}

新页表本身也需要物理页。只要处理器还在 Stage 1 的旧 CR3 下，内核就不能解引用
64 MiB 以上新分配的页表帧；因此建表阶段的分配上限必须暂时保持在恒等映射窗口。
待新根同时包含低端启动映射、内核映像、高半区直映、早期堆、保护测试页和
设备映射，并且所有步骤都成功后，内核才能写 CR3。

\begin{pdfdiagram}
\node[os state, text width=3.0cm] (facts)
  {验证输入事实\\E820、地址宽度\\managed limit};
\node[os memory, text width=3.0cm, right=13mm of facts] (meta)
  {低端选择元数据\\初始化全 PFN 账本\\保留启动对象};
\node[os compute, text width=3.0cm, right=13mm of meta] (build)
  {旧 CR3 下建表\\页表帧限于低端\\建立两类映射};
\node[os control, text width=3.0cm, below=16mm of meta] (commit)
  {提交新 CR3\\回读根与保护位\\切换页表访问策略};
\node[os state, text width=3.0cm, right=13mm of commit] (validate)
  {直映后自检\\查询映射、写读高端页\\释放并核对统计};
\draw[os strong bus] (facts) -- (meta);
\draw[os strong bus] (meta) -- (build);
\draw[os strong bus] (build) |- (commit);
\draw[os strong bus] (commit) -- (validate);
\draw[os feedback] (build.south) to[bend left=14]
  node[os label, right] {任一步失败均不采用半成品根} (facts.south);
\end{pdfdiagram}
\diagramnote{新 CR3 是启动内存事务的提交点。提交前所有页表页必须能经旧恒等
映射访问；提交后统一经 direct map 访问。失败路径不能让一部分高端映射成为
当前机器状态。}

在 64 GiB 主规格上，目标机还要从 4 GiB 以上分配一个页帧，查询对应直映叶项，
确认它是 supervisor、writable、NX、普通缓存，然后通过高半区地址写入并读回
两个 64 位模式，最后释放页帧。这个测试同时证明“高端 PFN 已进入分配器”“页表
确实采用了该映射”“处理器实际能够读写”“释放后所有权返回”四件事；只打印
managed bytes 不能替代它。当前主规格要求高端自检物理地址不低于
\code{0x0000000100001000}，并以 32768 个 2 MiB 叶项覆盖对齐的 RAM 内部区间；
边界与洞由 4 KiB 叶项和 not-present 区间分别表达。

\subsection{为什么 buddy 的 order 是地址结构，不是任意编号}

单页扫描只能回答“有没有一页可用”，无法有效回答“有没有一段对齐的连续页”。
当前工作树把一个连续块写成 \code{PhysicalFrameBlock\{physical\_address,
order\}}。若页大小为 4 KiB，则 order \(k\) 的含义完整展开为
\[
  \text{frame count}=2^k,\qquad
  \text{block bytes}=2^{k+12},\qquad
  \text{start PFN}\bmod 2^k=0.
\]
最后一项不是为了让地址好看。只有块头按自身大小对齐，同阶块才拥有唯一伙伴：
\[
  \text{buddy PFN}=\text{block PFN}\mathbin{\mathrm{xor}}2^k.
\]
例如 order 3 表示八页、32 KiB，起始 PFN 的低三位必须为零；异或
\(2^3\) 只翻转其上的伙伴选择位，低三位仍为零。若允许任意长度或任意起点，
“相邻空闲区”可能属于另一个父块，释放时便不能仅凭地址确定能否合并。

order 变化时没有页面字节被搬到新地址。分裂只是把一个高阶空闲块的位图事实改成
一个继续分裂的半块和一个放回空闲集合的半块；合并则做相反的元数据变化。块覆盖的
物理地址始终不变。只有调用者随后通过 direct map 读写页面，页面内容才真正变化。

\begin{longtable}{|L{0.20\linewidth}|L{0.25\linewidth}|L{0.29\linewidth}|L{0.16\linewidth}|}
\hline
\rowcolor{BookBlueTint}
\textbf{记录} & \textbf{它回答的问题} & \textbf{为什么不能由另一份记录代替} &
\textbf{错误时的风险} \\\tablerowrule
\endhead
两位 PFN 状态 & 每一页是 Unavailable、Free、Allocated 还是 Reserved &
逐页状态不知道一次申请原本覆盖多少页，也不知道其 order &
同一页被重复交付或保留页进入分配 \\\tablerowrule
各 order 的 Free 位图 & 哪些完整、对齐块当前可以直接取得或继续分裂 &
只扫描 Free 页不能证明这些页构成同一个可合并父块 &
跨洞、跨保留区或跨父块分配 \\\tablerowrule
各 order 的 Allocated 位图 & 活动块从哪个精确块头、以哪个 order 取得 &
只看一串 Allocated 页无法区分一个八页申请和八个单页申请 &
错阶释放把仍活动的子区间重新标为空闲 \\\tablerowrule
分裂、合并与活动统计 & 生命周期是否守恒，失败是否保持旧状态 &
计数是审计证据，不是所有权事实，不能反向代替位图 &
统计看似相等，块拓扑已经重叠 \\\hline
\end{longtable}

\begin{pdfdiagram}[scale=0.96]
  \node[os memory, text width=31mm] (ledger) at (-50mm,18mm)
    {两位 PFN 账本\\先识别洞、可用页\\与启动占用};
  \node[os control, text width=31mm] (freeze) at (0,18mm)
    {提交全部保留区\\元数据也先保留\\随后冻结修改};
  \node[os state, text width=31mm] (bitmaps) at (50mm,18mm)
    {按 order 建位图\\Free 与 Allocated\\解释块身份};
  \node[os compute, text width=31mm] (allocate) at (50mm,-18mm)
    {申请连续块\\预演路径后分裂\\只提交目标半块};
  \node[os control, text width=31mm] (release) at (0,-18mm)
    {按原地址与 order 释放\\验证后逐阶合并\\重复释放被拒绝};
  \node[os state, text width=31mm] (validate) at (-50mm,-18mm)
    {\code{ValidateBuddy}\\复算父子不重叠\\页数与统计守恒};
  \draw[os strong bus] (ledger.east) -- node[os label, above] {初始化} (freeze.west);
  \draw[os strong bus] (freeze.east) -- node[os label, above] {派生} (bitmaps.west);
  \draw[os strong bus] (bitmaps.south) -- (allocate.north);
  \draw[os strong bus] (allocate.west) -- node[os label, above] {归还} (release.east);
  \draw[os strong bus] (release.west) -- node[os label, above] {校验} (validate.east);
  \draw[os feedback] (validate.north) -- node[os label, left] {逐页核对} (ledger.south);
\end{pdfdiagram}
\diagramnote{buddy 位图是两位 PFN 账本在“连续块阶次”上的投影。初始化必须先完成
所有保留，再一次性建立投影；运行时分裂、分配、释放和合并同时更新两层状态，
完整校验器再从位图复算页数并检查父子块不重叠。}

\subsection{为什么元数据先保留，buddy 初始化后不再追加保留区}

页帧状态存储和 buddy 位图都占用真实 RAM。内核先计算两者尺寸，在低端恒等映射
可访问区中寻找一段连续空间，再把前半段交给两位 PFN 账本、后半段交给 buddy
位图。接着初始化逐页账本，保留固件、Stage 1、内核、BootInfo、启动栈、临时
缓冲和这段元数据自身；只有所有保留成功后，才从最终 Free 页生成各阶空闲块。

这个顺序解释了两个看似严格的拒绝。第一，buddy 初始化前若已经存在普通分配，
仅凭 Allocated 页无法恢复它原来的块头和 order，因此
\code{InitializeBuddy} 返回“已有分配”而不是猜测。第二，buddy 初始化后，
\code{ReserveRange} 会返回“初始化后禁止保留”；否则逐页账本已经把某页改成
Reserved，各阶 Free 位图却可能仍把包含它的父块列为空闲。未来可以实现“先拆除
覆盖该区间的空闲块，再事务化保留”的接口，但当前阶段选择冻结保留集合，使两层
状态只有一条可审计的初始化路径。

对 \(N\) 个受管 PFN，每一阶分别保存 Free 与 Allocated 两族位图，原始尺寸为
\[
  M(N)=2\sum_{k=0}^{\lfloor\log_2N\rfloor}
  \left\lceil\frac{\left\lceil N/2^k\right\rceil}{8}\right\rceil
  \text{ bytes}.
\]
前面的 2 来自两族位图，不是“留一倍保险”。仅有 Free 位图无法识别一次释放
携带的 order 是否等于原申请；Allocated 位图正是错阶、错位和重复释放的证据。
纯 64 GiB、4 KiB 页的容量测试得到 \(8\,\mathrm{MiB}+4\) 字节原始尺寸；
主 QEMU 规格的 PFN 下标还覆盖 3--4 GiB 地址洞后的 65 GiB 物理端点，因此
页对齐后的目标机输出为 \code{BUDDY\_STORAGE\_BYTES=0x821000}。地址洞不会
变成 Free，但仍占据按 PFN 编号的位图位置。

代码把通用上限写成 order 51，也不是假定机器拥有 \(2^{51}\) 页。一个 64 位
半开物理地址范围最多表达 \(2^{64}\) 字节；4 KiB 页已经占去 12 位，因此
order 52 会代表 \(2^{52}\) 页、恰好 \(2^{64}\) 字节，其独占终点无法装入
\code{uint64\_t}。order 51 是仍可安全移位并表达完整块范围的最大阶。运行时还会
取 \(\lfloor\log_2N\rfloor\) 与 51 的较小者；当前 65 GiB 地址跨度只得到
\code{BUDDY\_MAX\_ORDER=0x18}，即 order 24，而不是为不存在的更高阶建立块。

\subsection{为什么目标自检申请 order 3，又为什么地址不是常量}

目标机自检申请 order 3 的八页块，并分别通过 direct map 写读首尾页。单页测试
无法证明“中间七个页边界都属于同一连续块”，过大的高阶测试又会人为提高最低内存
要求；八页、32 KiB 是一个能触发真实分裂与合并、同时保持启动成本有界的选择。
释放以后，自检要求 Free、Allocated、Reserved 页数恢复，活动块数恢复，成功
申请与释放各增加一次，并让 \code{ValidateBuddy} 全量通过。

范围下界是 \code{0x0000000100001000}，即 4 GiB 再加一页。加这一页是为了证明
结果严格位于 32 位地址上界之外，而不是刚好落在边界；分配器还会把候选起点向上
调整到 32 KiB 块对齐。当前 64 GiB 主规格实际得到
\code{BUDDY\_SELF\_TEST\_ADDRESS=0x0000001000000000}。这个地址由 E820 空闲区、
启动保留、此前页表分配、范围和 order 共同筛选，不是代码固定的设备地址或 ABI。
换一份内存图仍可返回其他合规地址，测试只要求整个八页块落在半开范围内且保持
order 3 对齐。

宿主单元测试验证 64 页模型的拆分、合并、错阶、错位、重复释放和失败原子性；
生命周期测试加入 E820 地址洞、保留区与高地址半开范围；固定种子随机测试在
1024 页模型中执行 \(100000\) 次 order 0--7 混合申请释放，每 257 步全量复算
位图不变量。三项新增测试把完整套件从 \(77\) 项推进到 \(80\) 项；本次完整构建、
ELF/固件审计和真实 QEMU 路径均通过。

\subsection{从单调指针到可回收字节块}

物理直映只解决“内核怎样解引用一个已经映射的物理页”，并不解决页上对象何时
死亡。v0.6 的堆只保存下一个空闲地址：每次分配向高地址推进，失败时返回容量
不足，却没有合法的释放操作。这样的单调堆适合与内核同寿命的启动对象；若
Thread、FileTable 或动态内核栈沿用它，对象退出只会留下无法复用的字节。

当前 \code{KernelHeap} 仍使用原有 64 KiB RW/NX 高半区，但每个物理块增加
48 字节块头，保存自身长度、前块长度、请求长度、状态签名和双向空闲链接。
自身长度与前块长度构成边界标记，使释放路径能够直接找到相邻块；空闲块另按
地址递增进入双向链表。分配先遍历完整空闲集合，选择能容纳请求的最小块，再
计算对齐前缀与剩余后缀；所有边界和溢出检查通过以后，才一次性提交分裂结果。

\begin{longtable}{|L{0.18\linewidth}|L{0.28\linewidth}|L{0.29\linewidth}|L{0.15\linewidth}|}
\hline
\rowcolor{BookBlueTint}
\textbf{操作} & \textbf{提交前必须证明} & \textbf{提交后改变的事实} &
\textbf{失败时保持} \\\tablerowrule
\endhead
初始化 & 区间非空，基址、长度和最小块满足 16 字节与 64 字节约束 &
形成一个覆盖全部容量的空闲块 & 原对象仍未初始化 \\\tablerowrule
分配 & 请求非零、对齐是二次幂，候选前缀与后缀均合法且无溢出 &
至多产生前空闲块、活动块和后空闲块，并更新活动与峰值统计 &
输出指针与旧拓扑不变 \\\tablerowrule
释放 & 指针是活动负载的精确首地址，块头、前后边界和统计下界一致 &
活动块变为空闲块，并与前后空闲邻居形成唯一合并结果 &
非法或重复释放不改堆 \\\tablerowrule
校验 & 物理块链可完整遍历，空闲链有序、反向链接一致且无环 &
不修改状态，只报告整堆是否保持不变量 &
损坏状态显式可见 \\\hline
\end{longtable}

目标机先完成 16 字节与 4 KiB 对齐的真实写回，再逆序释放两个对象并要求活动
数与当前占用归零。宿主侧单元测试覆盖内部指针、重复释放、三向合并、地址复用
与统计；64 KiB 集成测试在乱序释放后重新申请最大连续负载；固定种子模型执行
\(100000\) 次随机申请与释放，逐步核对对齐、首尾数据和拓扑。堆增量落地时，
三项新增测试使套件由 v1.0 的 73 项增长到 76 项；命名门禁随后推进到 77 项，
本节前述 buddy 三项测试又把当前完整套件推进到 80 项。

当前实现仍不是无限增长、并发安全的通用内核分配入口。后备区固定为 64 KiB，
best-fit 线性扫描空闲链，调用者还必须保证启动路径串行；IRQ、NMI、panic 与
多 Thread 分配都不在当前契约内。当前 buddy 已负责页粒度连续块，但尚无
并发、zone、迁移类型与回收水位；未来 type cache 负责固定尺寸热对象，KVA
负责虚拟区间。通用堆不能替代这些层次。

\subsection{三项内存基础怎样继续长成通用内核分配}

\begin{longtable}{|L{0.21\linewidth}|L{0.24\linewidth}|L{0.27\linewidth}|L{0.18\linewidth}|}
\toprule
项目中的 v1.1 对象 & 已经解除的 v1.0 限制 & 正文继续引入的机制 & 始终保留的不变量 \\
\midrule
\endhead
动态两位 PFN 账本 & 固定状态数组只覆盖低 64 MiB & NUMA 节点、zone、水位与
内存热插拔 & 一个页帧同一时刻只有一个所有权状态 \\\tablerowrule
连续页块与两族位图 & 单页扫描不能稳定取得连续物理块 &
迁移类型、compaction、CMA 与并发空闲表 &
一个活动块只能按原始地址和 order 释放一次 \\\tablerowrule
\code{AllocateInRange} & 启动建表与高端普通分配不能区分 & DMA zone、低端保留池、
迁移类型与连续分配 & 约束范围外的页不能作为成功结果 \\\tablerowrule
高半区物理直映 & 高端物理页不能被内核 C++ 访问 & 临时映射、页缓存页、COW
复制与统一页表遍历 & 映射存在不等于页帧可分配 \\\tablerowrule
可回收 \code{KernelHeap} & 单调堆无法表达对象销毁、复用与碎片合并 &
type cache、动态内核栈、对象引用与失败回滚 &
释放只接受活动负载的精确首地址 \\\tablerowrule
混合 2 MiB/4 KiB 直映 & 全部小页导致建表成本随容量快速增长 & 大页拆分、
内存热插拔与更细权限 & 大页优化不能跨越地址洞或权限边界 \\
\bottomrule
\end{longtable}

因此，direct map 不是 \code{kmalloc}，buddy 也不是字节对象分配器，边界标记堆
仍不是 type cache。当前四个 PCB、内核栈和描述符容量尚未动态化，堆后备区也
不会按需增长。v1.1 的三项已接受增量分别解除“内核只能
看见低 64 MiB”“堆对象只能分配、不能归还”和“只能逐页扫描、不能恢复连续块”
三条限制，为后续动态对象表、COW 与按需分页建立可逐层验收的资源通路。
